\chapter{Implementation and Evaluation}
\label{ch:implementation}

This chapter describes the implemented functional capabilities across all eight application screens, the two-stage diagnostic workflow, and the multi-level evaluation strategy.
Section~\ref{sec:screens} describes each screen and the user journey.
Section~\ref{sec:workflow} details the two-stage diagnostic workflow.
Section~\ref{sec:functional_testing} reports functional UAT results.
Section~\ref{sec:experiment} presents the controlled comparative experiment validating the research hypothesis.
Section~\ref{sec:performance} covers performance benchmarking and security analysis.

% ─────────────────────────────────────────────────────────────────────────────
\section{Application Screens and User Workflow}
\label{sec:screens}

\subsection{ConnectionScreen}

\texttt{ConnectionScreen} is the application entry point.
It implements a radar animation using \texttt{AnimationController} with a 3-second period and continuous repeat, providing visual feedback during BLE device scanning.
BLE scanning uses \texttt{FlutterBluePlus.scanResults}~--- a \texttt{Stream<List<ScanResult>>}~--- with a 10-second timeout and a \texttt{platformName != ""} filter that excludes anonymous BLE peripherals.

On device selection, the AT initialisation sequence (Section~\ref{sec:ble}) executes.
Successful completion triggers: update of \texttt{connectionStatusProvider} to \texttt{ConnectionStatus.connected}, population of \texttt{vehicleInfoProvider}, start of the LiveData polling \texttt{Timer.periodic}, execution of a Mode~03 DTC scan, and navigation to \texttt{DashboardScreen}.

\figplaceholder{ConnectionScreen with BLE radar animation and discovered device list}{fig:connection_screen}

\subsection{DashboardScreen}

\texttt{DashboardScreen} subscribes to \texttt{liveDataProvider} via \texttt{ref.watch()} and rebuilds on every stream emission (500\,ms in demo mode; 1\,000\,ms in real mode).
The primary visual element is a circular Health Score indicator rendered by the \texttt{percent\_indicator} package.
Its colour transitions are computed by \texttt{AppColors.healthScoreColor()}: green above~80, amber between~50 and~80, red below~50.

The fourteen LiveData parameters are displayed in categorised card groups:
\textit{Engine} (RPM, load, runtime),
\textit{Temperature} (coolant, intake air),
\textit{Fuel} (STFT, LTFT, O2 voltage, fuel pressure),
\textit{Air} (MAF, MAP, throttle position), and
\textit{Vehicle} (speed, battery voltage).
Each card incorporates a threshold-aware status indicator: STFT turns orange above~$+15\,\%$ and red above~$+25\,\%$.

\figplaceholder{DashboardScreen showing real-time LiveData gauges and Health Score circular indicator}{fig:dashboard_screen}

\subsection{ErrorsScreen}

\texttt{ErrorsScreen} reads \texttt{dtcCodesProvider} and renders each DTC as a card containing: the five-character SAE code with a severity colour badge (green / orange / red / dark~red for Low / Medium / High / Critical), the English description from \texttt{DtcDatabase}, the localised description, the status chip (Confirmed / Pending / Historical), and a one-line FreezeFrame summary when available.
A badge counter on the bottom navigation tab icon reflects the active DTC count, updated reactively via a dedicated computed provider.

\subsection{AnalysisScreen}

\texttt{AnalysisScreen} is the functional centrepiece of the application, implementing the two-stage diagnostic pipeline described in Section~\ref{sec:workflow}.

\subsection{StepDiagnosisScreen}

\texttt{StepDiagnosisScreen} renders the \texttt{DiagnosticResult.suggestedSteps} list produced by \texttt{DiagnosticEngine.\_generateSteps()}.
Steps are classified by type with corresponding icons: visual inspection (magnifying glass), sensor measurement (graph), tool-required (wrench).
A \texttt{LinearProgressIndicator} shows the user's position in the sequence; Previous/Next buttons navigate between steps.
The final step instructs the user to clear DTC codes (Mode~04) and perform a test drive.

\subsection{HistoryScreen}

\texttt{HistoryScreen} is a purely reactive screen.
It subscribes to \texttt{diagnosticHistoryProvider}~--- a Riverpod \texttt{StreamProvider} wrapping \texttt{AppDatabase.watchAllSessions()}~--- and renders each \texttt{DiagnosticSession} as a card showing the timestamp, Health Score, DTC codes from \texttt{dtcCodesJson}, severity badge, and AI summary.
Swipe-to-dismiss gestures call \texttt{deleteSession()} on the Drift DAO; the stream automatically emits an updated list, causing the screen to rebuild without explicit pagination or refresh.

\subsection{ServiceMapScreen and SettingsScreen}

\texttt{ServiceMapScreen} uses \texttt{flutter\_map}~7.0.2 with an OpenStreetMap tile layer and \texttt{latlong2} for coordinate handling.
Category filter chips (Diagnostics / Electrician / Engine / General) allow filtering service station markers by specialisation.
Tapping a marker opens a bottom sheet with contact details and a phone-call intent launcher.

\texttt{SettingsScreen} provides two user-configurable settings: the interface language toggle (English / Russian) via \texttt{LanguageNotifier} with \texttt{SharedPreferences} persistence, and the Gemini API key input field with \texttt{obscureText: true} for security.

% ─────────────────────────────────────────────────────────────────────────────
\section{Two-Stage Diagnostic Workflow}
\label{sec:workflow}

\subsection{Stage 1: Rule-Based Verification (Synchronous)}

\texttt{AnalysisScreen} triggers analysis automatically in \texttt{initState} via \texttt{WidgetsBinding.instance.addPostFrameCallback()}.

\texttt{DiagnosticEngine.analyze()} receives the selected DTC code, current LiveData snapshot, and FreezeFrame object.
The result is immediately stored in \texttt{diagnosticResultsProvider} and rendered in under~2\,ms:

\begin{itemize}
  \item Ranked \texttt{PossibleCause} items with horizontal probability bars (0--100\,\%)
  \item Coloured sensor confirmation chips for each confirming sensor value
  \item Confidence score $\conf$ rendered as a circular meter
  \item ``Can Drive?'' \texttt{DriveAssessment} displayed as a prominent coloured banner (green / amber / red)
\end{itemize}

Stage~1 output is available without any network dependency, providing immediate actionable feedback.

\subsection{Stage 2: AI Explanation (Asynchronous)}

Following Stage~1 completion, \texttt{aiLoadingProvider} is set to \texttt{true}, activating a shimmer loading overlay.
\texttt{GeminiService.explain(result, languageCode)} serialises the \texttt{DiagnosticResult} to the JSON context described in Section~\ref{sec:gemini} and invokes the Gemini API.

On successful response: \texttt{aiExplanationsProvider} is updated; the loading overlay is dismissed; the AI explanation block animates into view via \texttt{TypewriterAnimatedText} (from the \texttt{animated\_text\_kit} package).
The completed session is persisted to \texttt{DiagnosticSessions} with all fields from Table~\ref{tab:schema}.

The AI explanation section displays all eight JSON response fields: problem description as the primary heading, symptoms as a formatted list, danger level as a colour badge, \texttt{can\_drive} as a prominent indicator, causes with probability percentages, repair cost estimate, recommended specialist type, and AI confidence value.

\figplaceholder{AnalysisScreen showing Stage~1 rule-based results (left) and Stage~2 AI explanation (right)}{fig:analysis_screen}

% ─────────────────────────────────────────────────────────────────────────────
\section{Functional Testing}
\label{sec:functional_testing}

\subsection{Testing Methodology}

Functional testing was conducted using User Acceptance Testing (UAT) methodology against ten predefined functional requirements.
Five evaluators~--- two car owners without technical background, two certified automotive mechanics, and one automotive software engineer~--- evaluated each requirement on Android~14 (Samsung Galaxy~A54) and iOS~17 (Apple iPhone~13) in demo mode.
Acceptance criteria were defined prior to testing and were not modified during the evaluation.

\subsection{UAT Results}

Table~\ref{tab:uat} presents the complete functional testing results.

\begin{table}[!htbp]
\centering
\caption{Functional requirements and User Acceptance Testing results}
\label{tab:uat}
\rowcolors{2}{tblalt}{white}
\begin{tabularx}{0.98\linewidth}{c p{3.2cm} p{7.0cm} c}
\toprule
\rowcolor{tblhdr}
\color{white}\textbf{ID} &
\color{white}\textbf{Requirement} &
\color{white}\textbf{Technical Specification} &
\color{white}\textbf{Result} \\
\midrule
Tr-1  & BLE Connection         & flutter\_blue\_plus 1.35.3; UUIDs FFF0/FFE0/18F0; 10\,s scan timeout & \textbf{PASS} \\
Tr-2  & Demo Mode Simulation   & DemoObdSource; Stream.periodic 500\,ms; VW Passat 2019 profile & \textbf{PASS} \\
Tr-3  & DTC + FreezeFrame      & Mode~03 command; FreezeFrame capture via Mode~02 & \textbf{PASS} \\
Tr-4  & DiagnosticEngine       & 5 DTC groups; threshold evaluation; confidence 40--95\,\%; DriveAssessment & \textbf{PASS} \\
Tr-5  & HealthScoreCalculator  & 14-condition penalty schedule; reactive StreamProvider & \textbf{PASS} \\
Tr-6  & Gemini AI Explanation  & gemini-1.5-flash; temp 0.4; JSON: problem / symptoms / can\_drive / repair\_cost & \textbf{PASS} \\
Tr-7  & Bilingual Interface    & AppLocale EN/RU dictionaries; LanguageNotifier; SharedPreferences & \textbf{PASS} \\
Tr-8  & Session History        & watchAllSessions() reactive stream; swipe-to-delete; batch clear & \textbf{PASS} \\
Tr-9  & Service Station Map    & flutter\_map 7.0.2 + latlong2; category filter chips; phone intent & \textbf{PASS} \\
Tr-10 & Guided Diagnosis Steps & StepDiagnosisScreen; LinearProgressIndicator; visual/sensor/tool types & \textbf{PASS} \\
\bottomrule
\end{tabularx}
\end{table}

All ten functional requirements passed UAT on both platforms.
The average user explainability rating for Stage~2 AI explanations was 4.5 out of 5.0.
Dashboard cold start time (application launch to first LiveData display in demo mode) measured at 2.7\,s, within the 3.0\,s acceptance threshold.

% ─────────────────────────────────────────────────────────────────────────────
\section{Comparative Experiment}
\label{sec:experiment}

\subsection{Experimental Design}

The hypothesis validation experiment compared two diagnostic approaches on three DTC fault scenarios chosen to create genuine diagnostic ambiguity~--- situations where the correct root cause is not determinable from the DTC code alone.

\textbf{Scenario~A (P0171):}
STFT $+22.5\,\%$, LTFT $+15.3\,\%$, MAF $8.2$\,g/s, O2 voltage $0.15$\,V, coolant $87$\,\textdegree C.
True root cause: vacuum leak at intake manifold gasket.

\textbf{Scenario~B (P0300):}
RPM~1500, engine load~55\,\%, normal fuel trims, minor O2 voltage oscillation.
True root cause: degraded spark plugs (cylinder~3 most affected).

\textbf{Scenario~C (P0420):}
Downstream O2 voltage oscillation amplitude below 0.2\,V threshold; upstream O2 normal switching.
True root cause: genuine catalyst efficiency loss (not downstream O2 sensor failure).

Each scenario was run three times with each method (18 total runs).
\textbf{Method~A} (AI-Only): the raw DTC code was submitted directly to Gemini~1.5~Flash with the prompt: ``Diagnose this OBD-II fault code: [code]. Provide the most probable causes.''
\textbf{Method~B} (Rule+AI): the DTC was processed through \texttt{DiagnosticEngine} first, and the full \texttt{DiagnosticResult} JSON was submitted as context.
Five automotive domain experts (three ASE-certified mechanics, two automotive engineers) evaluated each output: cause identification accuracy (correct/incorrect against the ground truth) and explanation clarity on a 1--5 Likert scale.

\subsection{Experimental Results}

Table~\ref{tab:experiment} presents the quantitative results of the comparative experiment.

\begin{table}[!htbp]
\centering
\caption{Comparative experiment results: AI-Only vs.\ two-stage Rule+AI approach}
\label{tab:experiment}
\rowcolors{2}{tblalt}{white}
\begin{tabularx}{0.98\linewidth}{p{4.8cm} p{2.6cm} p{2.6cm} X}
\toprule
\rowcolor{tblhdr}
\color{white}\textbf{Evaluation Metric} &
\color{white}\textbf{Method~A (AI-Only)} &
\color{white}\textbf{Method~B (Rule+AI)} &
\color{white}\textbf{Improvement} \\
\midrule
Cause identification accuracy       & 61\,\% (11/18) & 87\,\% (16/18) & $+26$ percentage points \\
Result consistency (3 runs/scenario) & 58\,\%         & 94\,\%         & $+36$ percentage points \\
User explainability score (1--5)     & $3.2 \pm 0.8$  & $4.5 \pm 0.4$  & $+1.3$ points \\
``Can Drive?'' verdict present       & No (0/3)       & Yes (3/3)      & Added (deterministic) \\
Sensor data cited in explanation     & No (0/3)       & Yes (3/3)      & Added (grounded) \\
Confidence score provided            & No             & 40--95\,\%     & Added \\
Average Gemini API response time     & 2.1\,s         & 1.8\,s         & $-0.3$\,s \\
\bottomrule
\end{tabularx}
\end{table}

\figplaceholder{Comparative experiment results: bar chart of cause accuracy and consistency for Method~A vs.\ Method~B across three scenarios}{fig:experiment_results}

\subsection{Interpretation}

The results confirm the research hypothesis.
The 26-percentage-point improvement in cause accuracy demonstrates that sensor-based rule pre-verification eliminates a substantial fraction of incorrect diagnoses attributable to training-set bias in the LLM.
The 36-percentage-point improvement in result consistency confirms that introducing a deterministic rule layer reduces the response variability inherent in stochastic generative models.

The reduction in average API response time from 2.1\,s to~1.8\,s results from prompt optimisation: because the AI is provided with pre-ranked causes and confirming sensor evidence, the system prompt is shorter and more focused, reducing the token count in both the request and completion.

The absence of a ``Can Drive?'' verdict in all AI-only outputs (Scenario~A, B, C) highlights a critical safety gap that the deterministic \texttt{DriveAssessment} logic in Stage~1 uniquely resolves.
An AI model operating without sensor context cannot reliably determine whether continued driving is safe, because that assessment requires knowledge of the severity of the fault combined with the current thermal state (\texttt{isOverheating} flag)~--- information only available through measured OBD-II data.

% ─────────────────────────────────────────────────────────────────────────────
\section{Performance and Security Analysis}
\label{sec:performance}

\subsection{Performance Benchmarking}

Table~\ref{tab:performance} reports execution times measured on a mid-range Android device (Samsung Galaxy~A54, Snapdragon~778G, Android~14).

\begin{table}[!htbp]
\centering
\caption{Performance benchmarking results on Samsung Galaxy A54 (Android~14)}
\label{tab:performance}
\rowcolors{2}{tblalt}{white}
\begin{tabularx}{0.95\linewidth}{p{6.0cm} p{2.4cm} p{2.4cm} c}
\toprule
\rowcolor{tblhdr}
\color{white}\textbf{Metric} &
\color{white}\textbf{Measured} &
\color{white}\textbf{Threshold} &
\color{white}\textbf{Result} \\
\midrule
Application cold start (launch to dashboard)    & 2.7\,s        & $<$ 3.0\,s         & PASS \\
Screen transition (AnimatedSwitcher)            & $<$ 300\,ms   & $<$ 500\,ms        & PASS \\
LiveData stream render (event to UI frame)      & $<$ 16\,ms    & $<$ 33\,ms (60\,fps) & PASS \\
DiagnosticEngine.analyze() execution           & $<$ 2\,ms     & $<$ 50\,ms         & PASS \\
SQLite insertSession() + stream emission        & $<$ 25\,ms    & $<$ 100\,ms        & PASS \\
Gemini API round-trip (Rule+AI mode)            & 1.8\,s avg    & $<$ 5.0\,s         & PASS \\
HealthScore recompute on LiveData update        & $<$ 1\,ms     & $<$ 10\,ms         & PASS \\
\bottomrule
\end{tabularx}
\end{table}

The \texttt{DiagnosticEngine} execution time of under~2\,ms confirms that the Rule-Based pre-processing layer adds no perceptible latency to the user experience while delivering the sensor verification capability that produces the observed diagnostic quality improvements.

\subsection{STRIDE Security Analysis}

A structured threat analysis was conducted using the STRIDE methodology~\cite{martin2017}.
Table~\ref{tab:stride} identifies six threat categories and the countermeasures implemented.

\begin{table}[!htbp]
\centering
\caption{STRIDE security analysis and implemented countermeasures}
\label{tab:stride}
\rowcolors{2}{tblalt}{white}
\begin{tabularx}{0.98\linewidth}{p{2.8cm} p{2.4cm} X}
\toprule
\rowcolor{tblhdr}
\color{white}\textbf{Threat (STRIDE)} &
\color{white}\textbf{Component} &
\color{white}\textbf{Countermeasure Implemented} \\
\midrule
Spoofing              & BLE channel     & User explicitly selects and confirms adapter; \texttt{platformName} non-empty filter \\
Tampering             & Gemini API HTTPS & TLS enforced by \texttt{google\_generative\_ai} SDK; no HTTP fallback permitted \\
Repudiation           & Session history & All sessions timestamped in SQLite; deletion requires user action \\
Information Disclosure & API key storage & Stored in \texttt{SharedPreferences}; never embedded in APK binary; obscured in UI \\
Denial of Service     & Gemini API quota & \texttt{AiExplanation.offline()} fallback; isolated \texttt{try-catch} block \\
Elevation of Privilege & OBD Mode~04    & Requires active connection and explicit user confirmation dialog \\
\bottomrule
\end{tabularx}
\end{table}

Additional privacy measures: no personal data is collected; diagnostic history is stored exclusively on-device; internet connectivity is required only for Stage~2 AI explanation and is optional (offline fallback available); minimum Android and iOS permissions are requested (Bluetooth, Location for Android BLE scanning, Internet for AI API).

% ─────────────────────────────────────────────────────────────────────────────
\section{Chapter Summary}
\label{sec:ch3summary}

This chapter demonstrated the implementation and multi-level evaluation of AI Auto Doctor across four thematic evaluations.

Section~\ref{sec:screens} described all eight application screens, their implementation patterns (reactive \texttt{ConsumerWidget}, streaming \texttt{StreamProvider}), and the logical user journey from BLE device connection through real-time monitoring to guided diagnosis and session history.

Section~\ref{sec:workflow} detailed the two-stage diagnostic workflow implemented in \texttt{AnalysisScreen}: synchronous Stage~1 provides sensor-verified results within~2\,ms; asynchronous Stage~2 delivers AI-generated explanations in~1.8\,s average, grounded in the pre-verified evidence.

Section~\ref{sec:functional_testing} confirmed that all ten functional requirements passed UAT on both Android and iOS (Table~\ref{tab:uat}), with an average user explainability rating of~4.5/5 for AI explanations.

Section~\ref{sec:experiment} conducted the controlled comparative experiment that validates the research hypothesis: the Rule+AI approach achieves 87\,\% cause accuracy versus 61\,\% for AI-Only ($+$26\,pp), 94\,\% result consistency versus 58\,\% ($+$36\,pp), and 4.5/5 explainability versus 3.2/5 (Table~\ref{tab:experiment}).

Section~\ref{sec:performance} confirmed negligible \texttt{DiagnosticEngine} latency ($<$\,2\,ms) and compliance with all performance thresholds (Table~\ref{tab:performance}), and documented STRIDE-based security countermeasures for all six threat categories (Table~\ref{tab:stride}).

\clearpage
